\documentclass{beamer}
\usepackage[utf8]{inputenc}

\usetheme{Berkeley}
\usecolortheme{default}

\setbeamertemplate{footline}[frame number]
\beamertemplatenavigationsymbolsempty


\title{What Makes a Paper Work?}
\subtitle{Structure, Flow, and Sign-Posting}
\author[]{IB 514: Scientific Writing}
\date{}

\AtBeginSection[]
{
  \begin{frame}
    \frametitle{Overview}
    \tableofcontents[currentsection]
  \end{frame}
}


\begin{document}

\frame{\titlepage}

%------------------------------------------------

\begin{frame}{Session Objectives}

(Hopefully) you analyzed a well-written paper before class.

\vspace{0.5cm}

Today we will identify:

\begin{itemize}
\item What makes scientific papers \textbf{readable}
\item What makes them \textbf{engaging}
\item What might make them \textbf{well-cited}
\end{itemize}

\vspace{0.5cm}

Focus on the \textbf{writing}, not the science.

\end{frame}


%------------------------------------------------
\section{Individual summary}
\begin{frame}{Individual summary (5 minutes)}

Look back at your worksheet.
Write one sentence answering:

\vspace{0.3cm}

\textbf{What is the single biggest (non-scientific) reason this paper was enjoyable to read?}

\vspace{0.5cm}

Then highlight in the paper (or your worksheet)...

\begin{itemize}
\item one \textbf{writing technique}
\item one \textbf{structural feature}
\end{itemize}

\hspace{2cm}...that contributed to that experience.


\vspace{0.5cm}

Think about examples such as:

\begin{itemize}
\item Introduction structure
\item Topic sentences
\item Paragraph transitions
\item Evidence progression
\end{itemize}

\end{frame}

%------------------------------------------------
\section{Small-group discussion}
\begin{frame}{Small-group discussion (12 minutes)}

Form groups of 3–4.
Each person briefly explains:

\begin{enumerate}
\item One structural feature
\item One sign-posting technique
\item Why the paper was enjoyable to read (science-independent)
\end{enumerate}

\vspace{0.5cm}

Ask each other:

\begin{itemize}
\item Where in the paper does this occur?
\item What exactly did the authors do?
\end{itemize}

\end{frame}

%------------------------------------------------
\begin{frame}{Identify patterns across papers (10 minutes)}

As a group, identify \textbf{3–5 patterns} across your papers.

Discuss:

\begin{itemize}
\item Structural similarities
\item Sign-posting techniques
\item Narrative flow
\item Figure-text alignment
\item Writing style
\end{itemize}

\vspace{0.5cm}

Key question:

\textbf{Which of these features might make a paper more likely to be cited?}

\vspace{0.5cm}
Add the patterns your group finds to the class whiteboard.

\end{frame}


%------------------------------------------------
\section{Takeaways}
\begin{frame}{Takeaways}

Strong scientific writing depends on:

\begin{itemize}
\item clear, recognizable structure
\item logical, sign-posted flow
\end{itemize}

\vspace{0.4cm}

Reading strong (and weak) papers carefully for their writing rather than their scientific content
is one of the best ways to improve your own writing.

\end{frame}
%------------------------------------------------
\section{Implementation}
\begin{frame}{Implementation (remaining time)}

Apply these lessons to your own manuscript.

\begin{enumerate}
    \item  Choose \textbf{three techniques} from today's discussion.
    \item  Identify specific places in your manuscript to apply them.
    \item  Write down a short revision plan.
        \begin{itemize}
            \item Rewrite topic sentence of paragraph $x$ to state main result
            \item Add conceptual figure
            \item Move interpretation from Results to Discussion
        \end{itemize}
\end{enumerate}
\end{frame}

\end{document}